% ============================================================
%  Estatística — Apostila Premium | PratiqueJá
%  Engine: XeLaTeX
% ============================================================
\documentclass[12pt]{article}
\usepackage{../../pratiqueja}
\usepackage{tikz}

% \hlm — destaque de resultado em modo-math (cor sem bold)
\newcommand{\hlm}[1]{{\color{darkblue}#1}}

\setsubject{Estatística}

\begin{document}
\pagenumbering{arabic}
\setstretch{1.2}
\color{bodytext}

\heroheader{Estatística}%
           {Médias, Dispersão, Frequências e Gráficos}%
           {Matemática · Avançado}

\setlength{\columnsep}{18pt}
\setlength{\columnseprule}{0.5pt}
\def\columnseprulecolor{\color{exborder}}

\begin{multicols}{2}

%─────────────────────────────────────────────────────────────
\TW{Def}{badgepurple}{Conceitos Iniciais}{População, amostra e variáveis}

\regra{A \textbf{estatística} organiza e interpreta dados para apoiar
decisões. \textbf{População} = conjunto total; \textbf{amostra} =
subconjunto observado.}

\SH{Tipos de variável}
\obs{%
  \textbf{Quantitativa discreta:} inteiros contáveis (nº de filhos).\\[2pt]
  \textbf{Quantitativa contínua:} qualquer valor de um intervalo (altura).\\[2pt]
  \textbf{Qualitativa nominal:} categorias sem ordem (cor dos olhos).\\[2pt]
  \textbf{Qualitativa ordinal:} categorias ordenadas (escolaridade).%
}

\exemp{%
  ``nº de gols'' → \hlm{discreta}\\[2pt]
  ``tempo em segundos'' → \hlm{contínua}\\[2pt]
  ``satisfação: ruim$<$\,\ldots\,$<$ótimo'' → \hlm{qualitativa ordinal}%
}

%─────────────────────────────────────────────────────────────
\TW{Méd}{darkblue}{Tendência Central}{Média, mediana e moda}

\regra{Indicam um \textbf{valor representativo} do conjunto.}

\SH{Média simples e ponderada}
\formula{$\bar{x} = \dfrac{x_1 + x_2 + \cdots + x_n}{n}$\\[6pt]
$\bar{x}_p = \dfrac{x_1 w_1 + \cdots + x_n w_n}{w_1 + \cdots + w_n}$}

\SH{Mediana (Me) — valor central do rol}
\obs{$n$ \textbf{ímpar}: elemento da posição $\frac{n+1}{2}$.\\
$n$ \textbf{par}: média das posições $\frac{n}{2}$ e $\frac{n}{2}+1$.}

\SH{Moda (Mo) — valor mais frequente}
\obs{Pode ser \textbf{amodal} (sem moda), \textbf{unimodal},
\textbf{bimodal} etc.}

\exemp{%
  \textbf{Notas} $5,6,6,7,10$:\\
  $\bar{x} = \dfrac{34}{5} = \hlm{6{,}8}$\\
  $n=5$ (ímpar) → posição $3$ → Me $= \hlm{6}$ \quad Mo $= \hlm{6}$%
}
\exemp{%
  \textbf{Ponderada} $7(\times2),\,5(\times3),\,8(\times5)$:\\
  $\bar{x}_p = \dfrac{14+15+40}{10} = \hlm{6{,}9}$%
}
\nota{A média sofre com \textbf{outliers}; a mediana é mais robusta;
a moda indica o valor mais típico.}

%─────────────────────────────────────────────────────────────
\TW{Dis}{badgegreen}{Medidas de Dispersão}{Amplitude, variância e desvio padrão}

\regra{Quantificam o quanto os dados se \textbf{afastam} do centro
(duas amostras de mesma média podem diferir muito).}

\formula{$A = x_{\max} - x_{\min}$}
\formula{$\sigma^2 = \dfrac{\sum (x_i - \bar{x})^2}{n} \qquad \sigma = \sqrt{\sigma^2}$}
\obs{O desvio padrão $\sigma$ tem a \textbf{mesma unidade} dos dados.}

\exemp{%
  \textbf{Notas} $4,6,7,8,10$: \ $\bar{x} = \dfrac{35}{5} = 7$\\
  $A = 10 - 4 = 6$\\
  desvios: $-3,\,-1,\,0,\,+1,\,+3$\\
  $\sigma^2 = \dfrac{9+1+0+1+9}{5} = \dfrac{20}{5} = 4$\\
  $\sigma = \sqrt{4} = \hlm{2}$%
}
\nota{$\sigma$ pequeno → dados concentrados; $\sigma$ grande →
dados espalhados.}

%─────────────────────────────────────────────────────────────
\TW{Fr}{badgepurple}{Distribuição de Frequências}{Absoluta, relativa e acumulada}

\regra{Muitos dados são organizados em \textbf{classes} $[a,b)$, que
agrupam os valores com $a \leq x < b$.}

\obs{$f_i$ = freq.\ \textbf{absoluta} (qtd.\ na classe);
$fr_i = f_i/n$ = \textbf{relativa};
$F_i = f_1 + \cdots + f_i$ = \textbf{acumulada}.}

\begin{center}\vspace{2pt}
\setlength{\tabcolsep}{8pt}\renewcommand{\arraystretch}{1.25}\footnotesize
\begin{tabular}{cccc}
\rowcolor{darkblue}
\textcolor{white}{\textbf{Classe}} & \textcolor{white}{$\mathbf{f_i}$} &
\textcolor{white}{$\mathbf{fr_i}$} & \textcolor{white}{$\mathbf{F_i}$}\\
$[0,2)$  & 1 & 0,10 & 1\\
\rowcolor{exbg}
$[2,4)$  & 2 & 0,20 & 3\\
$[4,6)$  & 4 & 0,40 & 7\\
\rowcolor{exbg}
$[6,8)$  & 2 & 0,20 & 9\\
$[8,10]$ & 1 & 0,10 & 10\\
\rowcolor{rulebg}
\textbf{Total} & \textbf{10} & \textbf{1,00} & —\\
\end{tabular}
\end{center}

\nota{$F_3 = 7$: sete alunos com nota abaixo de $6$. $fr_3 = 0{,}40$:
$40\%$ dos alunos na classe $[4,6)$.}

%─────────────────────────────────────────────────────────────
\TW{Grá}{darkblue}{Gráficos Estatísticos}{Histograma, polígono e boxplot}

\SH{Histograma e polígono de frequências}
\obs{\textbf{Histograma:} barras \emph{justapostas}, uma por classe.
\textbf{Polígono:} liga os pontos médios dos topos das barras.}

\begin{center}\vspace{2pt}
\begin{tikzpicture}[x=0.42cm,y=0.5cm]
\foreach \i/\h in {0/1,1/2,2/4,3/2,4/1}{
  \fill[accentblue!28!white] (\i*2,0) rectangle (\i*2+2,\h);
  \draw[darkblue,line width=0.8pt] (\i*2,0) rectangle (\i*2+2,\h);
}
\draw[line width=1pt] (0,0)--(10.4,0);
\draw[line width=1pt] (0,0)--(0,4.7);
\foreach \x in {0,2,4,6,8,10} \node[font=\scriptsize] at (\x,-0.5){\x};
\foreach \y in {1,2,3,4} \node[font=\scriptsize] at (-0.7,\y){\y};
\draw[vered,line width=1pt] (1,1)--(3,2)--(5,4)--(7,2)--(9,1);
\foreach \p in {(1,1),(3,2),(5,4),(7,2),(9,1)} \fill[vered] \p circle (2pt);
\end{tikzpicture}
\end{center}

\SH{Boxplot (diagrama de caixa)}
\obs{Resume a distribuição em \textbf{cinco valores}: mínimo, $Q_1$,
mediana $(Q_2)$, $Q_3$ e máximo. A caixa vai de $Q_1$ a $Q_3$; os
\textbf{bigodes} vão ao mín.\ e ao máx.\\[2pt]
Ex.: $2,4,5,6,7,8,9$ → inferior $\{2,4,5\}\Rightarrow Q_1{=}4$;
superior $\{7,8,9\}\Rightarrow Q_3{=}8$.}

\begin{center}\vspace{2pt}
\begin{tikzpicture}[x=0.5cm,y=0.5cm]
\draw[line width=0.8pt] (0,-0.15)--(10.3,-0.15);
\foreach \x in {0,2,4,6,8,10} {\draw (\x,-0.15)--(\x,-0.35); \node[font=\scriptsize] at (\x,-0.7){\x};}
\draw[darkblue,line width=1pt] (2,0.7)--(4,0.7);
\draw[darkblue,line width=1pt] (8,0.7)--(9,0.7);
\draw[darkblue,line width=1pt] (2,0.4)--(2,1.0);
\draw[darkblue,line width=1pt] (9,0.4)--(9,1.0);
\fill[accentblue!25!white] (4,0.2) rectangle (8,1.2);
\draw[darkblue,line width=1pt] (4,0.2) rectangle (8,1.2);
\draw[vered,line width=1.4pt] (6,0.2)--(6,1.2);
\node[font=\scriptsize] at (2,1.55){Mín};
\node[font=\scriptsize,darkblue] at (4,1.55){$Q_1$};
\node[font=\scriptsize,vered] at (6,1.55){$Q_2$};
\node[font=\scriptsize,darkblue] at (8,1.55){$Q_3$};
\node[font=\scriptsize] at (9,1.55){Máx};
\end{tikzpicture}
\end{center}

\nota{\textbf{IQR} $= Q_3 - Q_1$. Valores fora de
$[\,Q_1 - 1{,}5\,\text{IQR},\ Q_3 + 1{,}5\,\text{IQR}\,]$ são
\textbf{outliers}.}

\end{multicols}

\end{document}
